\section{Numerical results and discussion}
\label{sec:results}

The evaluation first separates the effects of mission-cost modeling,
spatial refinement and trajectory generation through controlled comparisons
within the same planning framework. A 72-case experiment then examines
mission-preference sensitivity, from which three representative preferences
are selected to illustrate the observed time--exposure--fuel trade-offs.
Search-stage estimates and final trajectory metrics are distinguished
throughout the analysis.

%% 当前紧凑结构：4.1 设置；4.2 三组对比；4.3 四张敏感性核心图；4.4 三个代表。
%% 4.2 的 A/B/C 分别支撑三个贡献，不再为每项贡献另立 subsection。
%% 执行顺序：验收数据 -> 4.1 -> 4.3 -> 4.4自动选择 -> 4.2必要的main对照。
%% B/C共用同一次main；4.4默认直接用sweep中三组最终缓存，不额外跑三次main。
%% 本节英文是比较协议和分析依据；实际结果段落等待实测数据，不预填提升结论。
%% 图均为 %% 占位，不会尝试加载不存在的文件。已合格的72组无需因改稿重跑。

\subsection{Simulation setup and evaluation protocols}
\label{sec:setup}

The simulation uses a $100\times100~\mathrm{km}$ horizontal workspace
containing five ground-based radar sites and five associated cylindrical
exclusion zones. The radar layout is generated with seed 42 and held fixed
throughout the preference comparison. Each radar has a hemispherical
coverage radius of $15~\mathrm{km}$, while its colocated hard zone has a
horizontal radius of $6~\mathrm{km}$ and a height of $5~\mathrm{km}$.
The power-law exposure model uses an inner-radius ratio of 0.3 and an
exponent of 4, with overlapping fields combined by their maximum.
These settings define a controlled synthetic environment; they are not
calibrated measurements of a particular radar or aircraft system.

The mission endpoints are $(2,2,0)$ and $(98,98,0)~\mathrm{km}$.
The en-route altitude interval is $1$--$10~\mathrm{km}$, with a
$5~\mathrm{km}$ horizontal transition distance at each endpoint.
The zero endpoint altitude is the reference used for the airborne
trajectory model, as described in Section~\ref{pr:sec:constraints}.
The vehicle and energy parameters are listed in Table~\ref{tab:aircraft};
the numerical planning settings are summarized in
Table~\ref{tab:planning}.

\begin{table}[t]
 \centering
 \caption{Parameters of the reduced fixed-wing model.}
 \label{tab:aircraft}
 \small
 \begin{tabular}{ll}
 \hline
 Parameter & Value \\
 \hline
 Constant mass & $1100~\mathrm{kg}$ \\
 Wing reference area & $11.5~\mathrm{m^2}$ \\
 Gravitational acceleration & $9.81~\mathrm{m\,s^{-2}}$ \\
 Stall / cruise / aircraft maximum speed & $33 / 55 / 83~\mathrm{m\,s^{-1}}$ \\
 Maximum ascent / descent rate & $12 / 15~\mathrm{m\,s^{-1}}$ \\
 Tangential acceleration / deceleration bound & $3 / 4~\mathrm{m\,s^{-2}}$ \\
 Maximum bank / flight-path angle & $30^\circ / 15^\circ$ \\
 Structural turning-radius floor & $50~\mathrm{m}$ \\
 Zero-lift drag coefficient / aspect ratio & $0.025 / 18$ \\
 Oswald efficiency & $0.85$ \\
 Climb / overall propulsive efficiency & $0.80 / 0.30$ \\
 Descent energy recovery fraction & $0$ \\
 Fuel lower heating value & $43~\mathrm{MJ\,kg^{-1}}$ \\
 Sea-level density / density scale height & $1.225~\mathrm{kg\,m^{-3}} / 8500~\mathrm{m}$ \\
 \hline
 \end{tabular}
\end{table}

\begin{table}[t]
 \centering
 \caption{Search and trajectory-evaluation settings. The sweep and the
 complete two-stage method use different input resolutions and locality limits.}
 \label{tab:planning}
 \small
 \begin{tabular}{ll}
 \hline
 Parameter & Value \\
 \hline
 Coarse horizontal / vertical spacing & $2000 / 500~\mathrm{m}$ \\
 Fine horizontal / vertical spacing & $200 / 50~\mathrm{m}$ \\
 Three-dimensional fine-search corridor radius & $600~\mathrm{m}$ \\
 Spline locality limit, fine input & $200~\mathrm{m}$ \\
 Spline locality limit, coarse sweep input & $3015.90~\mathrm{m}$ \\
 Search edge probes / quadrature subsegments & $10 / 9$ \\
 A* heuristic inflation factor & $1.0$ \\
 Initial smoothing budgets & $21$ \\
 Spline geometric validation samples & $4096$ \\
 Arc-length inversion probes & $8000$ \\
 Final target sampling distance & $100~\mathrm{m}$ \\
 Endpoint / terminal-heading tolerance & $2~\mathrm{m} / 12^\circ$ \\
 Initial speed & $33~\mathrm{m\,s^{-1}}$ \\
 Terminal speed & No prescribed equality \\
 \hline
 \end{tabular}
\end{table}

Two protocols address different experimental questions. The complete
two-stage protocol comprises coarse search, corridor-restricted fine search,
spline refinement and strict speed validation. It supports the controlled
comparisons in Section~\ref{sec:ablation}. The sensitivity protocol
uses coarse search followed by refinement and speed validation for each
of the 72 weight configurations specified in
Section~\ref{pr:sec:objective}. These results support both the sensitivity
analysis and the representative-preference comparison in
Sections~\ref{sec:sensitivity} and~\ref{sec:representative}. Its locality limit is scaled to the coarse
cell diagonal,
\begin{equation}
 r_s^{\mathrm{sweep}}=1.05\sqrt{2(2000)^2+(500)^2}
                    \simeq3015.90~\mathrm{m},
 \label{ex:eq:sweeptrust}
\end{equation}
whereas the complete method uses $r_s=200~\mathrm{m}$ relative to the fine
polyline. The two protocols share the threat and aircraft models, objective
scales, candidate-ranking rule and final metric definitions. Their
different search domains and locality limits are retained explicitly
when interpreting the results.

Search-stage and final-trajectory quantities are distinguished throughout
the analysis. Flight time $T$, exposure $C$ and fuel $F$ are reported in
seconds, exposure-seconds and kilograms, respectively; path length is
reported in kilometres. The exposure integral measures time-weighted
presence in the prescribed field, not a mission-level detection probability.
Cross-preference comparisons use the physical vector $(T,C,F)$: a value
of $J$ computed with one weight vector cannot directly rank a trajectory
scored with another. For a controlled comparison at fixed weights, the
shared anchor normalization provides the common scalar objective.

%% 4.1 操作与核对：
%% A. 环境来自 radar_environment_5zones.py::generate(seed=42)，正文使用 main.py
%%    实际端点 (2,2,0)/(98,98,0) km，不用环境类内部默认的 (5,5)/(95,95)。
%% B. 从运行缓存/原始环境导出五站坐标；归档精确值，不用截图反读坐标。
%% C. 填真实实验机器 CPU、RAM、操作系统、Python/NumPy/SciPy 版本。
%%    本次写作机器的软件版本不能自动充当先前实验的版本。
%% D. 归档 baseline.pkl 中的三个 anchor 原始 T/C/F 与六个尺度参考值；
%%    三条 anchor 未做最终样条处理，不混进“最终轨迹”比较表。
%% E. 计时应分别报告 anchor（一次性）、粗搜索、细搜索、样条、速度/评价。
%%    不把绘图/Excel 导出时间混进搜索时间，不把缓存加载时间当重新计算时间。
%%    只有实际补做重复测量后才能写重复次数、均值/标准差；同图重复计时
%%    反映计算波动，不是多环境统计证据。未测步骤写未记录，不填 0。
%% F. preflight_check.py 可作运行前检查；程序测试不能代替实验数据验收。
%% G. 建议最小数值核查：对选定轨迹另做更密集几何/积分采样，比较最小净空、
%%    最大曲率、最大角度和 T/C/F 的变化；固定几何时无需重跑 A* 或 72 组。
%%    100 -> 50 m 若重新分配速度，属于“离散化+速度策略”联合变化，须单列；
%%    不把它说成仅积分误差。未完成上述核查前不宣称数值收敛或连续安全认证。

%% FIGURE PLACEHOLDER: environment overview; let LaTeX assign figure numbers.
%% Panels: (a) 3-D hemispheres and hard cylinders with mission endpoints;
%% (b) exposure slices at stated altitudes with shared [0,1] color scale.
%% Hard-zone disk only appears when the slice altitude is <= 5 km.
%% Use physical axis units and distinguish hard zones from traversable exposure.
%% \begin{figure}[t]
%% \centering
%% \includegraphics[width=\linewidth]{figures/threat_environment.pdf}
%% \caption{Synthetic mission environment with five radar fields and colocated
%% hard exclusion zones. Soft exposure remains traversable and contributes
%% to the time-integrated objective.}
%% \label{fig:environment}
%% \end{figure}

\subsection{Overall planning performance and controlled comparisons}
\label{sec:ablation}

Three comparisons address distinct parts of the planning chain. The first
examines the effect of the mission objective, the second compares coarse
and fine raw routes from one run, and the third assesses the transition
from that fine route to a continuous trajectory with an assigned speed
profile. The latter two comparisons reuse intermediate outputs rather
than requiring separate planning experiments. These are controlled objective
and stage comparisons, rather than a complete factorial ablation: each
comparison is interpreted at the level of the factors it actually changes.

%% 本节按 A/B/C 三组写，每组遵循：问题 -> 共同条件 -> 数据 -> 实测现象 -> 解释。
%% 不以一张混合 raw/final 的表计算“总提升”，三组对比的指标口径分别明确。
%% 三锚点属于目标建模基础，不单列第四项创新；无须另加外部算法竞赛章节。

\subsubsection{Time-oriented A* versus proposed planning}
\label{sec:objective_comparison}

The reference is a time-oriented A* planner with weights $(1,0,0)$.
Hard-zone admissibility, altitude and climb constraints are retained,
whereas continuous soft exposure and fuel do not enter its search cost.
The proposed formulation uses the normalized composite cost in
Eq.~\eqref{pr:eq:objective}. Thus, the comparison concerns the information
encoded in the objective, rather than a limitation of A* to binary costs.
The radar exposure field remains traversable for both formulations and
is evaluated using the same model after planning.

Both formulations use the same environment, aircraft parameters, lattice
resolutions and refinement and speed-assignment procedures. Their coarse
raw paths can first be compared on the common coarse graph. Final mission
performance is then compared after each route has completed the two-stage
pipeline and passed its feasibility checks. Each route defines its own
fine corridor, so this final comparison measures the end-to-end effect
of the objective rather than holding the fine search domain fixed.

The principal measures are path length $L$, flight time $T$, accumulated
exposure $C$ and fuel consumption $F$. The analysis concerns their
trade-off, not simultaneous improvement in all four quantities. A lower
exposure accompanied by additional time or fuel would indicate the cost
of the selected mission preference, if observed in the data. A scalar
comparison of $J$ requires both trajectories to be evaluated under one
stated weight vector and the same anchor references.

%% A-1：基线是time-oriented A*，不是“A*只懂撞与不撞”；保留所有共同硬约束。
%%      不把整片软雷达区变成基线禁飞区，否则同时改变了可行空间。
%% A-2：多目标权重预先声明；可用4.4规则选出的平衡权重，再执行完整main。
%%      默认(0.3,1,0.2)只能称默认/示例偏好，不自动等于Balanced。
%% A-3：每次main先归档配置、baseline、coarse/final pkl和日志，避免固定文件名覆盖。
%%      原始coarse的T/C/F是search估计；正文主表比较两个通过验证的final结果。
%% A-4：同时加入C和F时只能归因于复合代价；不单独声称两个项各自带来多少收益。
%% A-5：表列Method、weights、L(km)、T(s)、C(exposure-s)、F(kg)，可附总计算时间。
%%      用100*(new-base)/base报告有符号变化；base=0时给绝对差，不加epsilon造百分数。
%% A-6：图用相同场景叠加轨迹，可做俯视/三维/高度多面板；风险切片标明高度。
%%      C_proposed<C_baseline且T_proposed>T_baseline只是可能的权衡，不能提前写成结果。
%%      失败记录保留失败阶段，不能用零值或raw路径代替成功final结果。
%% \begin{figure}[t]
%% \centering
%% \includegraphics[width=\linewidth]{figures/time_vs_risk_aware.pdf}
%% \caption{Time-oriented and mission-oriented final trajectories obtained
%% under the common two-stage planning and flight-envelope validation protocol.}
%% \label{fig:objective_comparison}
%% \end{figure}

\subsubsection{Coarse versus fine planning}
\label{sec:resolution_comparison}

The coarse raw route and the fine raw route are extracted from the same
main run. Both use identical mission weights and anchor scales. Their
comparison therefore concerns the local geometric changes produced by
fine-grid replanning within the three-dimensional corridor, before any
spline deformation or speed reassignment. Spatial overlays and altitude
profiles identify where the fine route departs from the coarse reference.

Search-stage $L$, $T$, $C$, $F$ and $J$ are reported together with the
stage-specific computation times and expansion counts. These observations
assess the practical effect and cost of local refinement; improvement in
every metric is not presumed. Without an unrestricted fine-grid reference,
the results do not establish a quantified speedup over full-domain fine
search or a globally optimal route topology.

%% B-1：使用同一次main的planned_path_coarse.pkl['path','stats']与
%%      planned_path.pkl['path_raw','stats_raw']，不是后者最终path/stats。
%% B-2：叠加粗路径、600m三维走廊和细原始路径，补局部放大/高度剖面。
%%      复用Section 3的coarse_corridor_fine图即可，不重复生成一张同内容图。
%% B-3：计算细原始点/边到粗折线的三维距离；二维投影不能代替该距离。
%%      最终样条可出粗走廊，因默认只强制相对fine raw的200m局部偏离。
%% B-4：表列Coarse raw/Fine raw的搜索估计和阶段开销；fine_time不是两阶段总耗时。
%%      coarse/fine并非简单插值，搜索图也非必然嵌套，故不能预写fine J必定下降。
%% B-5：可说限制了细搜索域并报告实测开销，不能无全域fine对照就声称快了X倍。

\subsubsection{Fine raw path versus B-spline trajectory}
\label{sec:refinement_comparison}

The fine raw path and its refined trajectory are obtained from the same
run. This comparison addresses the distinction between an admissible
piecewise-linear route and a continuous trajectory that passes the adopted
fixed-wing flight-envelope checks. Local geometry, endpoint and heading
retention, altitude and hard-zone screening are examined together with
minimum turning radius, analytic spline curvature and flight-path angle.
Curvature inferred from three raw vertices is a discretization diagnostic,
not continuous curvature at a polyline corner.

The final speed $v(s)$ and signed tangential acceleration $a(s)$ are
assessed against the stall, local curvature, vertical-rate and acceleration
limits. The initial speed is prescribed, while no terminal-speed equality
is imposed. This evaluation concerns numerical compliance with the reduced
flight-envelope model, not validation of a complete aircraft control system.
Final $T$, $C$ and $F$ are recomputed along the refined geometry and its
assigned speed profile. Their simultaneous reduction is not a requirement:
smoothing and retiming serve feasibility, and the geometric selection
score $J_s$ is distinct from the mission objective $J$.

The relation between geometry and exposure is examined along the final
trajectory. Since $dt=ds/v(s)$, a section with elevated threat intensity
contributes more exposure when its traversal speed is lower. Curvature-
or vertical-rate-induced speed reductions can therefore modify the
mission trade-off after route selection. Aligning the threat intensity,
curvature and speed profiles by arc length helps identify whether such
effects occur in the selected case. This interpretation concerns the
observed trajectory; it does not attribute the entire raw-to-final cost
change to a single processing stage.

%% C-1：同一planned_path.pkl取path_raw、path、smooth_result、curvature_smooth、
%%      speed_profile、accel_profile；无需为此再运行main。
%% C-2：图(a)原始/样条局部转弯，(b)解析曲率与1/R_ref，(c)速度与失速/局部上限，
%%      (d)带符号加速度；角度和垂直率可增加小面板或在约束表给极值与阈值。
%% C-3：重点R_min、kappa、gamma、v(s)、a(s)，再列端点误差<=2m、航向误差<=12deg、
%%      最大局部偏离<=200m、altitude余量、硬区检查。不能仅凭曲线平滑宣布可执行。
%% C-4：加速+3/减速-4m/s^2、上升12/下降15m/s分别检查；垂直率包括段内峰值。
%%      原始折线若画非严格速度曲线，明确标成diagnostic，不能充当可飞行基线。
%% C-5：raw的搜索估计与final的变速积分分列；两者差含几何/速度变化，不宣称纯平滑收益。
%%      报告候选预算、选中sigma和修复状态；没有消融则不宣称选择器已优于其他策略。
%% C-6：对至少一个最终案例加密几何/积分采样复核，无需重跑A*；若重新分配速度，
%%      单列为离散化与速度的联合变化。有限采样通过不等于连续安全证明。
%% C-7：复用现有v(s)/曲率图的横轴，标出高暴露弧段，必要时增加小面板Pd(s)。
%%      检查“弯曲/爬升限速 -> 局部dt增加 -> 暴露积分变化”是否真的出现在该案例。
%%      若高暴露区没有明显限速，也如实报告，不能为串联模块强行编造因果故事。
%%      由当前final缓存重积分即可，不需要额外场景或重跑72组。
%% \begin{figure}[t]
%% \centering
%% \includegraphics[width=\linewidth]{figures/refinement_feasibility.pdf}
%% \caption{Fine-route refinement and flight-envelope validation, including
%% analytic spline curvature and the speed and acceleration profiles.}
%% \label{fig:refinement}
%% \end{figure}

\subsection{Mission-preference sensitivity}
\label{sec:sensitivity}

The preference experiment fixes $w_C=1$ and varies
\begin{equation}
 w_T\in\{0,0.1,\ldots,1.1\},\qquad
 w_F\in\{0,0.1,\ldots,0.5\},
 \label{ex:eq:weights}
\end{equation}
resulting in $12\times6=72$ configurations. Each configuration follows
the coarse-search, spline-refinement and speed-assignment protocol in
Section~\ref{sec:setup}. The same environment, aircraft model, anchor
references and final evaluation rules are retained. In particular,
$(w_T,w_F)=(0,0)$ is an exposure-only objective because $w_C$ remains one.

Three heatmaps summarize the final time, exposure and fuel responses to
the two varied coefficients. A fourth plot presents the $T$--$C$ distribution
with fuel consumption encoded by color. Together, these views describe
the objective trade-offs without comparing scalar scores computed under
different preferences. Plateaus and abrupt changes, where present, should
be examined in relation to route selection and subsequent refinement;
monotonic changes in the final metrics are not assumed. The experiment
characterizes the sampled preferences in the stated environment, rather
than a global Pareto frontier or performance across independent scenarios.

%% 4.3只保留下面四张核心图，不再要求正文增加切片图、高度图或行为分区图。
%% 步骤1【验收】：analysis_data_loader核对72个唯一组合、成功状态、schema、平滑版本、
%% run_signature及当前baseline匹配；Excel“全部数据”与对应case缓存一致。
%% 缺失/失败先说明与处理，不能删行后声称72个有效结果，不混用旧工作簿。
%% 步骤2【三幅热力图】：x=w_T,y=w_F，全部12*6网格，单位T(s)、C(exposure-s)、F(kg)。
%% 同一指标的配套图使用一致色标；T/C/F不同量纲，不强行使用同一数值色标。
%% analyze_weight_sensitivity.py已有合并输出；如需三张独立图，可复用同一数据另行导出，
%% 本次只预留图位，并未修改绘图脚本或生成新图。不要从低分辨率截图拼接论文图。
%% 步骤3【散点】：横轴T(s)、纵轴C(exposure-s)、颜色F(kg)，保留所有配置的信息。
%% 重合点可用频数/标记说明；不将72组当独立随机重复实验，不自动添加回归显著性。
%% 步骤4【写作】：先报实测范围与极值位置；再固定一个权重沿行/列解释变化；
%% 最后用散点说明时间--暴露--燃油的代价。检查路线的分析可在后台做，不新增正文图。
%% 若计算非支配点，按三目标<=且至少一项<定义，并公开容差；不要从二维外观断言。
%% 本节的所有指标都是coarse+后处理final，不是72次完整fine流程。
%% \begin{figure}[t]
%% \centering
%% \includegraphics[width=\linewidth]{figures/time_heatmap.pdf}
%% \caption{Final flight time across the 72 configurations of the coarse-search protocol.}
%% \label{fig:time_heatmap}
%% \end{figure}
%% \begin{figure}[t]
%% \centering
%% \includegraphics[width=\linewidth]{figures/exposure_heatmap.pdf}
%% \caption{Time-integrated threat exposure across the 72 preference configurations.}
%% \label{fig:exposure_heatmap}
%% \end{figure}
%% \begin{figure}[t]
%% \centering
%% \includegraphics[width=\linewidth]{figures/fuel_heatmap.pdf}
%% \caption{Final fuel consumption across the 72 preference configurations.}
%% \label{fig:fuel_heatmap}
%% \end{figure}
%% \begin{figure}[t]
%% \centering
%% \includegraphics[width=\linewidth]{figures/time_exposure_scatter.pdf}
%% \caption{Time--exposure trade-offs in the sampled preference set, with fuel
%% consumption represented by color.}
%% \label{fig:tradeoffs}
%% \end{figure}

\subsection{Representative mission preferences}
\label{sec:representative}

Three selection rules identify risk-dominant, time-dominant and balanced
preferences from the same 72-case dataset. For case index $i\in\mathcal I$,
the first two rules select
\begin{equation}
 i_R\in\mathop{\arg\min}_{i\in\mathcal I} C_i,\qquad
 i_T\in\mathop{\arg\min}_{i\in\mathcal I} T_i.
 \label{ex:eq:extremecases}
\end{equation}
For the balanced rule, each physical objective $f_{ik}$,
$k\in\{T,C,F\}$, is scaled within the evaluated batch:
\begin{equation}
 f'_{ik}=\begin{cases}
 \displaystyle\frac{f_{ik}-\min_{j\in\mathcal I}f_{jk}}
 {\max_{j\in\mathcal I}f_{jk}-\min_{j\in\mathcal I}f_{jk}},
 &\max_j f_{jk}>\min_j f_{jk},\\
 0,&\text{otherwise}.
 \end{cases}
 \label{ex:eq:batchscale}
\end{equation}
The balanced case minimizes the Euclidean distance to the batch ideal point,
\begin{equation}
 D_i=\sqrt{(T'_i)^2+(C'_i)^2+(F'_i)^2},\qquad
 i_B\in\mathop{\arg\min}_{i\in\mathcal I}D_i.
 \label{ex:eq:balancedcase}
\end{equation}
This batch scaling is used only for representative selection and does
not replace the three-anchor normalization in the planner. Equal primary
values are resolved lexicographically by $(T,C,F,w_T,w_F)$ for the balanced
rule, $(T,F,w_T,w_F)$ for the exposure rule, and $(C,F,w_T,w_F)$ for the
time rule. These definitions match the implemented selection procedure.

The selected final trajectories are taken directly from the corresponding
sweep caches. One three-dimensional overlay and one table summarize their
geometry, weights and final $(T,C,F)$ values under that common protocol.
The labels describe selection rules within the sampled dataset, not
independent operating environments or globally optimal mission solutions.
If two rules select the same configuration, their coincidence is reported
rather than forcing distinct cases. The balanced choice expresses equal
importance in batch-scaled objective space; it is neither the minimizer
of $J$ across different weights nor a universally preferred mission setting.

%% 步骤1：运行select_representative_weights.py，取“绕飞/快穿/平衡”对应的三行。
%% 代码另外输出“省油”用于诊断，但本节只呈现Risk/Time/Balanced，不另加第四案例。
%% 使用自动规则，不手挑好看的权重；平台邻居数与指标指纹不参与选择。
%% 步骤2：读取对应sweep_cases_v5_3_1逐工况缓存，以TRAJECTORY_FILE选择输入；
%% 三条曲线与表中T/C/F均用同一coarse+后处理协议，不混入main结果。
%% 不把筛选表中四舍五入的值当重新积分输入；从已验收缓存/原始数据取精度。
%% 步骤3：叠加同一环境中的三条最终三维轨迹，统一视角、坐标、硬区与雷达图例。
%% 重合代表如实说明，可以保留三行但注明相同权重/轨迹，不声称三个独立解。
%% 步骤4：表仅Case,w_T,w_C,w_F,T,C,F，T用秒,C用exposure-s,F用kg。
%% 若坚持使用w_R表头，应在全文统一替换；当前w_C对应代码w_risk，避免符号混用。
%% 步骤5：文字用一段解释三种偏好的实际路线与代价，再以一段说明取舍和适用范围。
%% 不称Risk解一定全程绕飞、Time解一定穿越中心；“绕飞/快穿”是脚本名称而非几何定理。
%% 本节不额外跑main。需要完整两阶段验证的权重用于4.2，并单独归档。
%% 下表 -- 是尚未填入的占位，绝非零值；有真实结果后替换并移除此备注。
\begin{table}[t]
 \centering
 \caption{Representative preferences selected from the 72-case coarse-search
 experiment. Metrics refer to the final refined and retimed trajectories.}
 \label{tab:representatives}
 \small
 \begin{tabular}{lrrrrrr}
 \hline
 Case & $w_T$ & $w_C$ & $w_F$ & $T$ (s) & $C$ (exposure-s) & $F$ (kg) \\
 \hline
 Risk-dominant & -- & -- & -- & -- & -- & -- \\
 Time-dominant & -- & -- & -- & -- & -- & -- \\
 Balanced & -- & -- & -- & -- & -- & -- \\
 \hline
 \end{tabular}
\end{table}
%% \begin{figure}[t]
%% \centering
%% \includegraphics[width=\linewidth]{figures/representative_preferences.pdf}
%% \caption{Final three-dimensional trajectories for the risk-dominant,
%% time-dominant and balanced preferences selected from the coarse-search experiment.}
%% \label{fig:representatives}
%% \end{figure}

%% 第四章收束：围绕三个问题总结实测证据，保留效率/暴露/燃油的实际代价；
%% 不预写全面优越、跨场景鲁棒、真实飞行验证或连续安全保证。摘要/结论最后补数字。
